\section{Related Work}

\paragraph{Mechanistic Interpretability and Linear Representation Hypothesis.}
Mechanistic interpretability aims to understand neural networks by reverse-engineering their internal computations into human-interpretable components~\citep{olah2020zoom, elhage2021mathematical}.
Significant progress has been made in identifying interpretable circuits in transformers~\citep{elhage2021mathematical, olsson2022context} and in locating factual knowledge in their parameters~\citep{meng2022locating}.
A prominent assumption underlying this line of work is the \emph{linear representation hypothesis}~\citep{park2024lrh, elhage2022toy}, which posits that high-level concepts are encoded as linear directions in a model's activation space.
This view is supported by probing, editing, and steering studies that extract or manipulate concepts along learned directions in language models~\citep{hewitt2019structural, turner2024activation, nanda2023emergent, li2023inference}.

\paragraph{Sparse Autoencoders for Interpreting VLMs.}
SAEs have emerged as a scalable tool for interpreting activations in pre-trained models, with a line of work showing that SAEs recover mono-semantic features in language models~\citep{huben2023sparse, lieberum2024gemma, templeton2024scaling, gao2025scaling, markssparse}.
Building on these advances, SAEs have been extended to VLMs pre-trained with the contrastive objective~\citep{radford21clip}, where they have been used to extract interpretable per-modality concepts from image and text embeddings~\citep{rao2024discover, zaigrajew2025interpreting, pach2026sparse}.
A separate line of work shifts the focus from single-modality concept extraction to characterizing cross-modal structure---how concepts in image and text representations relate, align, or diverge~\citep{costa2026from, kaushik2026learning, dhimoila2026crossmodal}.

\paragraph{Internal Structure of Joint Embedding Space.}
Joint vision-language embeddings have been studied at both the embedding and feature levels.
At the embedding level, a prominent phenomenon is the \emph{modality gap}, where image and text embeddings occupy disjoint regions of the shared space~\citep{liang2022mind}.
Prior work attributes this gap either to \emph{modality-specific} features that are inactive in the other modality~\citep{jiang2023understanding, zhang2024connect, ramasinghe2024accept, schrodi2025two, qian2023intra}, or to misalignment of \emph{shared} features~\citep{eslami2025mitigate, yamaguchi2025post, grassucci2025closing, grassucci2026closing, grassucci2026the, mistretta2025cross}.

At the feature level, recent work analyzes representations using sparse autoencoders~\citep{zaigrajew2025interpreting, papadimitriou2025interpreting}.
These analyses reveal a related phenomenon, \emph{modality split}, where similar semantic concepts activate different latent units across modalities~\citep{kaushik2026learning, dhimoila2026crossmodal}.
Existing approaches aim to reduce this split by aligning latent activations, implicitly assuming that shared concepts correspond to common feature directions.
This assumption overlooks that shared concepts may have different feature directions across modalities.
In this work, we provide empirical and theoretical support for this view, and propose a method suited to this setting.
